%! AP Statistics — Unit 5, Lesson 5.4 — Group Activity (Answer Key)
\documentclass[10pt]{article}
\usepackage{apstats-article}
\usepackage{apstats-key}

\begin{document}

\pageheader{Unit 5, Lesson 5.4}{Group Activity --- Answer Key}
\namepartnerperiod

\vspace{0.1in}

\begin{tcolorbox}[colback=redbg, colframe=black!40, boxrule=0.6pt, arc=2mm,
    left=3mm, right=3mm, top=2mm, bottom=2mm, title=\textbf{Tier R --- Remediate}]
\begin{enumerate}
  \item \ans{Unbiased}
  \item \ans{Biased (low)}
  \item \ans{Biased (high)}
  \item \ans{$\mu_{\bar{x}} = 74$; $\sigma_{\bar{x}} = 12/\sqrt{9} = 4$.}
  \item \ans{$\sigma_{\bar{x}} = 12/\sqrt{36} = 2$. Spread decreased by a factor of $1/2$ (or: standard error halved when $n$ quadrupled).}
\end{enumerate}
\end{tcolorbox}

\vspace{0.1in}

\begin{tcolorbox}[colback=redbg, colframe=black!40, boxrule=0.6pt, arc=2mm,
    left=3mm, right=3mm, top=2mm, bottom=2mm, title=\textbf{Tier A --- Approaching Proficiency}]
\begin{enumerate}
  \item \ansline{$\bar{x}$ is an unbiased estimator of $\mu$. The sampling distribution of $\bar{x}$ is centered at $\mu = 62$ cm, meaning that on average, $\bar{x}$ neither over- nor under-estimates $\mu$.}
  \item \ansline{Biased --- biased high. The sample minimum averages 48 cm but the population minimum is 40 cm. On average, the sample minimum overestimates (is higher than) the population minimum.}
  \item \ans{$n=16$: $\sigma_{\bar{x}} = 8/\sqrt{16} = 2$ cm. $n=64$: $\sigma_{\bar{x}} = 8/\sqrt{64} = 1$ cm. With the larger sample, estimates cluster twice as tightly around the true mean --- the estimate is more precise.}
\end{enumerate}
\end{tcolorbox}

\vspace{0.1in}

\begin{tcolorbox}[colback=redbg, colframe=black!40, boxrule=0.6pt, arc=2mm,
    left=3mm, right=3mm, top=2mm, bottom=2mm, title=\textbf{Tier E --- Extension}]
\begin{enumerate}
  \item \ans{Estimator A is unbiased (centered at 100). Estimator B is biased low (centered at 97 $<$ 100).}
  \item \ansline{Not necessarily. Lower variability means estimates are more consistent (precise), but if they systematically miss the true value by 3 units, they are less accurate. In AP Statistics, unbiasedness is generally preferred; a biased estimator with lower variability might win on mean squared error, but the College Board emphasizes unbiasedness.}
  \item \ans{$z_1 = (95-100)/12 \approx -0.417$; $z_2 = (105-100)/12 \approx 0.417$; $P(-0.417<Z<0.417) \approx 0.6772 - 0.3228 = 0.3232 \approx 32.3\%$.}
  \item \ans{$z_1 = (95-97)/4 = -0.50$; $z_2 = (105-97)/4 = 2.00$; $P = 0.9772 - 0.3085 = 0.6687 \approx 66.9\%$. Estimator B actually has a higher probability of landing within 5 units of $\mu$ despite being biased --- illustrating the bias-variance trade-off.}
\end{enumerate}
\end{tcolorbox}

\end{document}
